%% Les trois mouvements de translation : rectiligne, circulaire, curviligne.
%% Dans chaque case, la même pièce en « L » (percée d'un trou) est dessinée
%% en position de départ (trait plein) et d'arrivée (pointillé) ; quatre points
%% de la pièce laissent quatre trajectoires colorées superposables.
\documentclass[tikz,border=4pt]{standalone}
\usepackage{liaisons}
\begin{document}
\begin{tikzpicture}[x=1cm,y=1cm, scale=0.78,
  traj/.style={line width=0.9pt, line cap=round},
  pt/.style={circle, inner sep=0pt, minimum size=3.2pt, fill}]

% --- la pièce en « L » (contour + trou, règle even odd) ---------------
\newcommand\piece[1]{%
  \path[even odd rule, line join=round, #1]
    (0,0) -- (1.15,0) -- (1.15,0.39) -- (0.39,0.39) -- (0.39,1.15) -- (0,1.15) -- cycle
    (0.20,0.20) circle (0.13);}
\tikzset{
  depart/.style={draw=black!65, fill=black!12, line width=0.8pt},
  arrivee/.style={draw=black!55, dashed, line width=0.7pt}}

% --- les quatre points suivis ----------------------------------------
\newcommand\pts{%
  \coordinate (Pa) at (0.20,1.13);
  \coordinate (Pb) at (1.12,0.20);
  \coordinate (Pc) at (0.03,0.03);
  \coordinate (Pd) at (0.72,0.36);}

%% ====================== 1. Translation rectiligne ===================
\begin{scope}[shift={(0,0)}]
  \node[align=center, font=\small] at (1.42,2.70) {Translation\\rectiligne};
  \pts
  \foreach \p/\c in {Pa/solideE, Pb/solideA, Pc/solideB, Pd/solideD}
    { \draw[traj, \c] (\p) -- ++(1.55,0.85); }
  \piece{arrivee, shift={(1.55,0.85)}}
  \piece{depart}
  \foreach \p/\c in {Pa/solideE, Pb/solideA, Pc/solideB, Pd/solideD}
    { \node[pt, \c] at (\p) {}; \node[pt, \c] at ($(\p)+(1.55,0.85)$) {}; }
\end{scope}

%% ====================== 2. Translation circulaire ===================
%% arcs de même rayon R, centres différents : parallèles, non concentriques
\begin{scope}[shift={(3.35,0)}]
  \node[align=center, font=\small] at (1.42,2.70) {Translation\\circulaire};
  \pts
  \def\R{2.2}
  \foreach \p/\c in {Pa/solideE, Pb/solideA, Pc/solideB, Pd/solideD}
    { \draw[traj, \c] (\p) arc (-90:-40:\R); }
  \piece{arrivee, shift={(1.686,0.786)}}
  \piece{depart}
  \foreach \p/\c in {Pa/solideE, Pb/solideA, Pc/solideB, Pd/solideD}
    { \node[pt, \c] at (\p) {}; \node[pt, \c] at ($(\p)+(1.686,0.786)$) {}; }
\end{scope}

%% ====================== 3. Translation curviligne ===================
\begin{scope}[shift={(6.70,0)}]
  \node[align=center, font=\small] at (1.42,2.70) {Translation\\curviligne};
  \pts
  \foreach \p/\c in {Pa/solideE, Pb/solideA, Pc/solideB, Pd/solideD}
    { \draw[traj, \c, shift={(\p)}]
        plot[domain=0:1.5, samples=45, smooth] (\x,{0.5*sin(\x*120)+0.33*\x}); }
  \piece{arrivee, shift={(1.5,0.495)}}
  \piece{depart}
  \foreach \p/\c in {Pa/solideE, Pb/solideA, Pc/solideB, Pd/solideD}
    { \node[pt, \c] at (\p) {}; \node[pt, \c] at ($(\p)+(1.5,0.495)$) {}; }
\end{scope}
\end{tikzpicture}
\end{document}
